\documentclass[../../../include/open-logic-section]{subfiles}

\begin{document}

\olfileid{sth}{ordinals}{iso}
\olsection{Isomorfisme Urutan}

Kanggo njlentrehake \emph{sepira} kukuhé urutan becik,
kita bakal miwiti kanthi cara mbandhingake urutan-urutan
becik.

\begin{defn}
\emph{Urutan becik} yaiku pasangan $\tuple{A, <}$ kang
$<$ ngurutake $A$ kanthi becik. Urutan becik
$\tuple{A, <}$ lan $\tuple{B, \lessdot}$ iku
\emph{isomorfis miturut urutan} yen lan mung yen ana
!!a{bijection} $f \colon A \to B$ kang nyukupi:
$x < y$ yen lan mung yen $f(x) \lessdot  f(y)$.
Ing kahanan iki kita nulis
$\ordeq{\tuple{A, <}}{\tuple{B, \lessdot}}$, lan
ngarani $f$ minangka \emph{isomorfisme urutan}.
\end{defn}
\noindent
Ing salajengipun, supaya ringkes, kita bakal nyebut
``isomorfisme'' tinimbang ``isomorfisme urutan''.
Kanthi intuitif, isomorfisme iku !!{bijection}
kang njaga struktur. Iki sawetara kasunyatan prasaja
ngenani isomorfisme.

\begin{lem}\ollabel{isoscompose}
Komposisi isomorfisme iku uga isomorfisme: yen
$f \colon A \to B$ lan $g \colon B \to C$ iku
isomorfisme, mula $(g \circ f) \colon A \to C$
uga isomorfisme.
\end{lem}
\begin{prob}
	Buktikna \olref[sth][ordinals][iso]{isoscompose}.
\end{prob}
\begin{proof}
Ditinggalake minangka gladhen.
\end{proof}
\begin{cor}\ollabel{ordisoisequiv}
	$\ordeq{X}{Y}$ iku relasi ekuivalensi.
\end{cor}
\begin{prop}\ollabel{ordisounique}
Yen $\tuple{A, <}$ lan $\tuple{B, \lessdot}$ iku
urutan becik kang isomorfis, mula isomorfisme ing
antarane loro iku tunggal.
\end{prop}

\begin{proof}
Upamane $f$ lan $g$ iku isomorfisme $A \to B$.
Kita bakal mbuktekake kanthi induksi, yaiku\ nggunakake
\olref[wo]{propwoinduction}. Pilihen $a\in A$ lan
upamane (kanggo induksi) yen $(\forall b < a)f(b) = g(b)$.
Pilihen $x \in B$.

Yen $x \lessdot f(a)$, mula $f^{-1}(x) < a$,
dadi $g(f^{-1}(x)) \lessdot g(a)$ amarga $f$ lan $g$
iku isomorfisme. Nanging amarga $f^{-1}(x) < a$,
miturut pangira-ira induksi kita
$x =f(f^{-1}(x)) = g(f^{-1}(x))$. Dadi
$x \lessdot g(a)$. Semono uga, yen $x \lessdot g(a)$,
mula $x \lessdot f(a)$.

Mula $(\forall x \in B)(x \lessdot f(a) \liff x \lessdot
g(a))$. Saka kene $f(a) = g(a)$ miturut
\olref[sfr][rel][ord]{prop:extensionality-strictlinearorders}.
Mula $(\forall a \in A)f(a) = g(a)$ miturut
\olref[wo]{propwoinduction}.
\end{proof}\noindent 
Iki menehi gambaran yen urutan becik iku kukuh.
Nanging kanggo nerangake luwih adoh, kita perlu
ngenalake notasi liyane.

\begin{defn}
Yen $\tuple{A, <}$ urutan becik kanthi $a \in A$,
upamane $A_a = \Setabs{x \in A}{x < a}$. Kita
ngarani $A_a$ \emph{perangan wiwitan} sejati saka
$A$ (lan $A$ dhewe uga dianggep perangan wiwitan
ora sejati saka $A$). Upamane $<_a$ iku watesan
$<$ marang perangan wiwitan iki, yaiku
$\funrestrictionto{\mathord{<}}{A_a^2}$.
\end{defn}
\noindent
Kanthi notasi iki, kita bisa negesake lan mbuktekake
yen ora ana urutan becik kang isomorfis karo
perangan wiwitan sejatine dhewe.

\begin{lem}\ollabel{wellordnotinitial}
Yen $\tuple{A, <}$ urutan becik kanthi $a \in A$,
mula $\ordneq{\tuple{A, <}}{\tuple{A_a, <_a}}$.
\end{lem}

\begin{proof}
Kanggo bukti kanthi kontradiksi, upamane
$f \colon A \to A_a$ iku isomorfisme. Amarga $f$
bijeksi lan $A_a \subsetneq A$, miturut
\olref[wo]{wo:strictorder} pilihen $b \in A$
minangka !!{element} paling cilik miturut $<$ ing
$A$ kang nyukupi $b \neq f(b)$. Kita bakal nuduhake
$(\forall x \in A)(x<b \liff x < f(b))$; mula miturut
\olref[sfr][rel][ord]{prop:extensionality-strictlinearorders}
kita entuk $b = f(b)$, sawijining kontradiksi.

Upamane $x < b$. Mula $x = f(x)$ miturut pilihan
$b$. Lan $f(x) < f(b)$ amarga $f$ isomorfisme.
Mula $x < f(b)$.

Upamane $x < f(b)$. Mula $f^{-1}(x) < b$
amarga $f$ isomorfisme; invers iki bisa dienggo amarga
perangan wiwitan ngemu kabeh unsur sadurunge
bayangan kasebut. Dadi $f^{-1}(x) = x$
miturut pilihan $b$. Mula $x < b$.
\end{proof}

Asil sabanjure nuduhake, kanthi ringkes, yen
``perangan wiwitan'' saka isomorfisme uga isomorfisme:

\begin{lem}\ollabel{wellordinitialsegment}
Upamane $\tuple{A, <}$ lan $\tuple{B, \lessdot}$
urutan becik. Yen $f \colon A \to B$ isomorfisme
lan $a \in A$, mula
$\funrestrictionto{f}{A_{a}} : A_a \to B_{f(a)}$
uga isomorfisme.
\end{lem}

\begin{proof}
Amarga $f$ isomorfisme:
	%Amarga $f$ isomorfisme, $b < a$ yen lan mung yen $f(b) \lessdot f(a)$, mula
\begin{align*}
	\funimage{f}{A_a} &= \funimage{f}{\Setabs{x \in A}{x < a}}\\
	&= \funimage{f}{\Setabs{f^{-1}(y) \in A}{f^{-1}(y) < a}} \\
	&= \Setabs{y \in B}{y \lessdot f(a)} \\
	&=B_{f(a)} 
\end{align*}
Lan $\funrestrictionto{f}{A_a}$ njaga urutan amarga
$f$ uga njaga urutan.
\end{proof}

Rong asil sabanjure nuduhake yen urutan becik mesthi
bisa dibandhingake:

\begin{lem}\ollabel{lemordsegments}
Upamane $\tuple{A, <}$ lan $\tuple{B, \lessdot}$
urutan becik. Yen
$\ordeq{\tuple{A_{a_1}, <_{a_1}}}{\tuple{B_{b_1}, \lessdot_{b_1}}}$
lan $\ordeq{\tuple{A_{{a_2}}, <_{a_2}}}{\tuple{B_{{b_2}},
\lessdot_{b_2}}}$, mula
${a_1}  < {a_2} \text{ yen lan mung yen }{b_1} \lessdot {b_2}$.
\end{lem}

\begin{proof}
Kita buktekake arah \emph{kiwa menyang tengen}; arah
kosok baline padha. Upamane
$\ordeq{\tuple{A_{a_1}, <_{a_1}}}{\tuple{B_{b_1},
\lessdot_{b_1}}}$ lan $\ordeq{\tuple{A_{{a_2}},
<_{a_2}}}{\tuple{B_{{b_2}}, \lessdot_{b_2}}}$,
kanthi $f \colon A_{{a_2}} \to B_{{b_2}}$ minangka
isomorfisme kita. Upamane ${a_1} < {a_2}$; mula
$\ordeq{\tuple{A_{a_1}, <_{a_1}}}{\tuple{B_{f({a_1})},
\lessdot_{f({a_1})}}}$ miturut
\olref{wellordinitialsegment}. Dadi
$\ordeq{\tuple{B_{b_1}, \lessdot_{b_1}}}{\tuple{B_{f({a_1})},
\lessdot_{f({a_1})}}}$, mula ${b_1} = f({a_1})$
miturut \olref{wellordnotinitial}. Saiki
${b_1} \lessdot {b_2}$ amarga jangkauan $f$ iku
$B_{{b_2}}$.
\end{proof}

\begin{thm}\ollabel{thm:woalwayscomparable}
Kanggo sembarang rong urutan becik, salah sijine
isomorfis karo perangan wiwitan (ora kudu sejati)
saka liyane.
\end{thm}

\begin{proof}
Upamane $\tuple{A, <}$ lan $\tuple{B, \lessdot}$
urutan becik. Kanthi Pamisahan, temtokna
\[
	f = \Setabs{\tuple{a, b} \in A \times B}{
		\ordeq{\tuple{A_a, <_a}}{\tuple{B_b, \lessdot_b}}}.
\]
Miturut \olref{lemordsegments}, $a_1 < a_2$ yen lan
mung yen $b_1 \lessdot b_2$ kanggo kabeh
$\tuple{a_1, b_1}, \tuple{a_2, b_2} \in f$.
Mula $f \colon \dom{f} \to \ran{f}$ isomorfisme.

Yen $a_2 \in \dom{f}$ lan $a_1 < a_2$, mula
$a_1 \in \dom{f}$ miturut
\olref{wellordinitialsegment}; dadi $\dom{f}$
perangan wiwitan saka $A$. Semono uga, $\ran{f}$
perangan wiwitan saka $B$. Kanggo kontradiksi,
upamane loro-lorone perangan wiwitan \emph{sejati}.
Pilihen $a$ minangka !!{element} paling cilik miturut
$<$ ing $A \setminus \dom{f}$, dadi
$\dom{f} = A_a$; lan pilihen $b$ minangka !!{element}
paling cilik miturut $\lessdot$ ing
$B \setminus \ran{f}$, dadi $\ran{f} = B_b$.
Mula $f \colon A_a \to B_b$ isomorfisme, saengga
$\tuple{a, b} \in f$, sawijining kontradiksi.
\end{proof}

\end{document}
